\documentclass[11pt]{article}

\usepackage[a4paper,margin=1in]{geometry}
\usepackage{amsmath}
\usepackage{amssymb}
\usepackage{booktabs}
\usepackage{enumitem}
\usepackage[hidelinks]{hyperref}

\emergencystretch=2em
\Urlmuskip=0mu plus 2mu

\title{Spaces, Duals, Morphisms, and \texttt{BlockTensor}}
\author{}
\date{\today}

\newcommand{\catC}{\mathcal{C}}
\newcommand{\Hom}{\operatorname{Hom}}
\newcommand{\id}{\operatorname{id}}
\newcommand{\dual}[1]{#1^{\vee}}
\newcommand{\Space}{\mathsf{Space}}
\newcommand{\DualSpace}{\mathsf{DualSpace}}
\newcommand{\Domain}{\mathsf{Domain}}
\newcommand{\Codomain}{\mathsf{Codomain}}
\newcommand{\MorphismSpec}{\mathsf{MorphismSpec}}
\newcommand{\BlockTensor}{\mathsf{BlockTensor}}
\newcommand{\SpaceId}{\mathsf{SpaceId}}
\newcommand{\BlockKey}{\mathsf{BlockKey}}
\newcommand{\CouplingDescriptor}{\mathsf{CouplingDescriptor}}

\begin{document}

\maketitle

\begin{abstract}
This note refines the local tensor model in
\emph{Planar Tensor Networks, Braiding, and Interface Invariants}.
The planar paper established that a symmetry-aware tensor is a morphism
with an ordered boundary and that global contraction semantics belong to
a checked planar diagram or contraction plan.  The refinement here
separates two distinctions that the earlier \(\Space/\mathsf{CoSpace}\)
spelling conflated: categorical object duality,
\(\Space/\DualSpace\), and morphism side,
\(\Domain/\Codomain\).  Moving a boundary factor from one side of a
morphism to the other is consequently an explicit wire-bending operation
that toggles object duality.  This model applies unchanged to the first
bosonic U(1) \(\BlockTensor\), where it reduces to signed charge
bookkeeping, and to later non-Abelian or braided implementations, where
it carries fusion, pivotal, and braid data.
\end{abstract}

\tableofcontents

\section{Relationship to the Planar-Network Design}

The earlier planar-network note supplies the global architecture:

\begin{enumerate}[nosep]
  \item a tensor stores a local ordered domain and codomain;
  \item a planar network stores the embedding and external boundary;
  \item a contraction plan inserts and lowers the required local
        compositions, recouplings, braids, and traces.
\end{enumerate}

Those conclusions remain unchanged.  This note corrects and completes
the local boundary vocabulary.  The earlier term
\(\mathsf{CoSpace}\) attempted to encode ``the other direction'' in one
type.  There are, however, two independent questions:

\begin{enumerate}[nosep]
  \item Is the categorical object \(X\) or its dual \(\dual{X}\)?
  \item Does this occurrence appear in the domain or codomain?
\end{enumerate}

The answers must not be collapsed.  A domain may contain either \(X\)
or \(\dual{X}\), and so may a codomain.  The durable vocabulary is
therefore \(\Space/\DualSpace\) for object duality and
\(\Domain/\Codomain\) for morphism side.

\section{Five Pieces of Boundary Data}

A symmetry-aware tensor boundary contains five different forms of
information.  Keeping them separate prevents structural metadata from
silently acquiring mathematical semantics.

\subsection{Symmetry or Tensor Category}

The symmetry context supplies the category \(\catC\).  It owns:

\begin{itemize}[nosep]
  \item sector fusion and the tensor unit;
  \item object duality;
  \item invariant-channel multiplicities;
  \item evaluation and coevaluation maps;
  \item eventually associators, pivotal data, twists, and braiding.
\end{itemize}

For bosonic U(1), most of this structure is trivial.  Fusion is charge
addition, the tensor unit has charge zero, duality negates charge, and
every allowed invariant channel has multiplicity one.  This is a
specialization of the interface, not a different interface.

\subsection{Structural Decomposition}

A decomposition records which sectors and degeneracy dimensions a
space contains.  The currently implemented Uni20 types provide useful
structural forms:

\begin{center}
\begin{tabular}{lll}
\toprule
Form & Structural data & Typical role \\
\midrule
\texttt{BlockSpace} & \((q,d_q)\) sectors & MPS virtual bond \\
\texttt{QNumList} & ordered local charges & physical site \\
\texttt{QNum} & one irrep label & irreducible operator \\
dense extent & no symmetry sectors & bundled dense axis \\
\bottomrule
\end{tabular}
\end{center}

Structural equality is not semantic identity.  Two MPS bonds can have
the same U(1) sectors and degeneracies while belonging to different
cuts of the chain.

\subsection{Identity-Bearing Space}

A \(\Space\) is an immutable semantic object over one decomposition.
It carries a \(\SpaceId\).

\begin{itemize}[nosep]
  \item Copying a space preserves its identity.
  \item Constructing an isomorphic space creates a distinct identity.
  \item Exact composition normally requires the same identity.
  \item Composition through distinct isomorphic spaces requires an
        explicit isomorphism.
  \item A display label is not the identity.
\end{itemize}

The implementation should expose the conceptual distinction
\[
  \operatorname{same\_space}(X,Y)
  \qquad\hbox{versus}\qquad
  \operatorname{isomorphic\_space}(X,Y).
\]
The first checks identity; the second checks compatible categorical and
structural data.

A process-local address is not a sufficient \(\SpaceId\).  MPI
construction and distributed decompositions require an identity that
can be serialized or deterministically reproduced on every
participating MPI rank.

\subsection{Object Duality}

For a space \(X\), \(\dual{X}\) is represented by
\(\DualSpace\langle X\rangle\).  It retains the identity of its primal
space:
\[
  \dual{(\dual{X})} \cong X,
  \qquad
  \operatorname{primal\_id}(\dual{X})
    = \operatorname{space\_id}(X).
\]

The dual is not a newly allocated, unrelated space whose decomposition
happens to contain dual sectors.  For U(1), its sector \(q\) is
represented by \(-q\).  In a richer category the symmetry context also
supplies the evaluation, coevaluation, and normalization data required
to use the dual.

\subsection{Morphism Side and Leg Occurrence}

A tensor is a morphism
\[
  A :
  X_0 \otimes \cdots \otimes X_{m-1}
  \longrightarrow
  Y_0 \otimes \cdots \otimes Y_{n-1}.
\]
Its local boundary is the ordered pair
\[
  \Domain[X_0,\ldots,X_{m-1}],
  \qquad
  \Codomain[Y_0,\ldots,Y_{n-1}].
\]

Every entry is one leg occurrence.  Several occurrences may refer to
the same space, just as an identity map has the same space in its
domain and codomain.  Occurrence identity, position, display label,
space identity, and object duality are distinct.

\section{The All-Out Boundary}

Fusion and block legality are most naturally expressed by orienting
every boundary strand outward.  Under this convention,
\[
  \Codomain[X] \longmapsto X,
  \qquad
  \Domain[X] \longmapsto \dual{X}.
\]
If an entry already names a dual object, the two dualities cancel.
Effective outward duality is therefore the exclusive-or of the
domain-side flag and the explicit dual-object flag.

For the morphism above, the all-out boundary is
\[
  \dual{X}_{m-1} \otimes \cdots \otimes \dual{X}_0
  \otimes
  Y_0 \otimes \cdots \otimes Y_{n-1},
\]
where the reversal of the dualized product records
\[
  \dual{(X_0 \otimes \cdots \otimes X_{m-1})}
  \cong
  \dual{X}_{m-1} \otimes \cdots \otimes \dual{X}_0.
\]
An implementation may use a different canonical cyclic starting point,
but it must choose and document one order and preserve it through every
operation.

A populated block corresponds to an invariant channel
\[
  \Hom_{\catC}
  \left(
    \mathbf{1},
    \dual{X}_{m-1}\otimes\cdots\otimes\dual{X}_0
    \otimes Y_0\otimes\cdots\otimes Y_{n-1}
  \right).
\]
The generic selection query is consequently a category operation over
oriented sectors, not a hard-coded arithmetic predicate.

\subsection{U(1) Selection Rule}

For bosonic U(1), the invariant-channel query reduces to
\[
  \sum_{j=0}^{n-1} q(Y_j)
  -
  \sum_{i=0}^{m-1} q(X_i)
  = 0.
\]
For an MPS site
\[
  A : L \otimes P \longrightarrow R,
\]
this is
\[
  q_R = q_L + q_P.
\]
For an MPO site
\[
  W : L \otimes K \longrightarrow R \otimes B,
\]
it becomes
\[
  q_L + q_K = q_R + q_B.
\]
These are instances of one oriented-boundary rule.

\section{Repartition Is an Explicit Bend}

Changing which side of a morphism contains a leg is not a relabel or a
bit flip.  It is a rigid-category isomorphism built from evaluation or
coevaluation:
\[
  \Hom_{\catC}(A\otimes X,B)
  \cong
  \Hom_{\catC}(A,B\otimes\dual{X}),
\]
\[
  \Hom_{\catC}(A,B\otimes X)
  \cong
  \Hom_{\catC}(A\otimes\dual{X},B).
\]

The operation is provisionally named \texttt{repartition}.  Moving an
occurrence across the boundary toggles its explicit
\(\Space/\DualSpace\) status so that the all-out boundary is preserved.

Moving several factors also exposes the contravariant ordering rule:
\[
  \dual{(X\otimes Y)}
  \cong
  \dual{Y}\otimes\dual{X}.
\]
The repartition API must therefore specify the resulting ordered
boundary, not only a new domain size.

In bosonic U(1), a bend can normally be implemented as checked metadata
and an axis-order view.  In a nontrivial pivotal category it may carry
normalization maps or twists.  It must always route through the
category hook even when that hook returns the identity.

\subsection{Operations That Must Remain Distinct}

\begin{itemize}
  \item \texttt{relabel}: changes diagnostic metadata only.
  \item \texttt{conj}: changes numerical value semantics while
        preserving the morphism boundary.
  \item \texttt{repartition}: bends selected strands and toggles
        object duality.
  \item \texttt{adjoint}: reverses the morphism and applies the
        category's adjoint and recoupling rules.
  \item \texttt{braid}: exchanges strands using the category's
        braiding.
\end{itemize}

These operations can have similarly shaped dense-array
implementations in the trivial case.  That does not make their
contracts interchangeable.

\section{Contraction as Tensor Product and Composition}

Full contraction is composition.  Given
\[
  A:X\longrightarrow Y,
  \qquad
  B:Y\longrightarrow Z,
\]
the result is
\[
  B\circ A:X\longrightarrow Z.
\]
Composition checks ordered boundary compatibility and space identity,
not only dimensions or sectors.

\subsection{Partial Contraction}

Partial contraction is tensor product with identities followed by
composition.  Suppose
\[
  A:L\longrightarrow M\otimes X,
  \qquad
  B:X\otimes R\longrightarrow N.
\]
Construct conceptually
\[
  A\otimes\id_R:
  L\otimes R
  \longrightarrow
  M\otimes X\otimes R,
\]
\[
  \id_M\otimes B:
  M\otimes X\otimes R
  \longrightarrow
  M\otimes N.
\]
Their composition is the requested contraction:
\[
  (\id_M\otimes B)\circ(A\otimes\id_R):
  L\otimes R
  \longrightarrow
  M\otimes N.
\]

Uni20 does not need to materialize either identity tensor.  The checked
planner proves the categorical expression, determines the external
boundary, and lowers directly to block worklists.

\subsection{General Planar Networks}

A planar network can be sliced into layers of tensor products followed
by compositions.  The general planner may additionally require:

\begin{itemize}[nosep]
  \item associators and \(F\)-moves;
  \item repartition bends;
  \item symmetric permutations or explicit \(R\)-moves;
  \item categorical traces for closed strands.
\end{itemize}

The local \(\BlockTensor\) owns its ordered morphism boundary.  The
network owns cyclic orders, global adjacency, and crossings.  A local
tensor cannot decide whether a proposed pairing is planar without that
network context.

\section{Why Bosonic Abelian Tensors Use the Full Model}

It is tempting to implement U(1) as a flat list of signed legs and add
the categorical model later.  That would make the first implementation
smaller only by postponing its central invariants.

The full model immediately provides:

\begin{itemize}[nosep]
  \item exact bond-space identity checks;
  \item explicit input and output boundaries;
  \item one contraction language for U(1) and richer categories;
  \item correct creation and sharing of decomposition spaces;
  \item a category-shaped selection-rule interface.
\end{itemize}

For bosonic U(1), it should compile to the inexpensive representation:

\begin{itemize}[nosep]
  \item multiplicity-free fusion;
  \item trivial associators and symmetric braiding;
  \item charge negation for object duality;
  \item an empty coupling descriptor;
  \item signed charge sums for block legality;
  \item metadata-only canonical bends and permutations.
\end{itemize}

A convenience contraction API may insert the unique canonical bend or
permutation in this symmetric setting.  It still produces a checked
plan.  A braided category must reject an expression that leaves a
crossing unspecified.

\section{Consequences for \texttt{BlockTensor}}

The block-sparse container should be parameterized by a morphism
boundary rather than a flat list of direction-prefixed leg types.
Illustrative C++ spelling is:

\begin{verbatim}
template <
    class T,
    class MorphismSpec,
    class Storage = HostOnly,
    class Coupling = Trivial>
class BlockTensor;

using MpsSiteSpec = MorphismSpec<
    Domain<BlockSpace, LocalSpace>,
    Codomain<BlockSpace>>;

using MpoSiteSpec = MorphismSpec<
    Domain<BlockSpace, LocalSpace>,
    Codomain<BlockSpace, LocalSpace>>;
\end{verbatim}

The exact class names remain design sketches.  The required split is:

\begin{itemize}[nosep]
  \item the static spec fixes order and decomposition kinds;
  \item tensor values carry the actual identity-bearing spaces;
  \item each boundary occurrence refers to one of those spaces or its
        dual;
  \item storage and placement policies do not change the boundary.
\end{itemize}

\subsection{Block Keys and Coupling Data}

A block key is an extensible structure
\[
  \BlockKey
  =
  (\text{sector index per boundary occurrence},
   \CouplingDescriptor).
\]
For multiplicity-free U(1), the coupling descriptor is empty.  For
non-Abelian fusion it identifies an invariant channel or fusion basis.
The container must not assume there is exactly one block for every
legal flat sector tuple.

Block construction calls a category-shaped query such as
\begin{verbatim}
symmetry.invariant_channels(oriented_sectors)
\end{verbatim}
and constructs only the channels it returns.

\subsection{Storage Orthogonality}

The same boundary and block key can be stored as individual host
allocations, offsets into a coalesced buffer, CUDA-resident blocks, or
MPI-distributed blocks.  A storage policy chooses ownership,
allocation, and placement.  It must preserve:

\begin{itemize}[nosep]
  \item every \(\SpaceId\);
  \item domain and codomain order;
  \item explicit object duality;
  \item sector indices and coupling descriptors;
  \item the global block identity used by distributed worklists.
\end{itemize}

Storage movement is not a symmetry transformation.

\section{Decompositions Create New Spaces}

An SVD is not only a dense matrix factorization.  It creates a new
intermediate space.  For
\[
  A:X\longrightarrow Y,
\]
one coherent convention is
\[
  V:X\longrightarrow B,
  \qquad
  S:B\longrightarrow B,
  \qquad
  U:B\longrightarrow Y,
\]
\[
  A=U\circ S\circ V.
\]

The factor naming may follow the final linalg API, but the boundary
contract is fixed:

\begin{enumerate}[nosep]
  \item \(B\) receives one new \(\SpaceId\).
  \item Every factor occurrence of the decomposition bond shares that
        identity.
  \item Truncation determines the structural decomposition of \(B\).
  \item An independently created isomorphic space is not silently
        substituted for \(B\).
\end{enumerate}

QR, LQ, polar decomposition, and DMRG bond truncation obey the same
rule.  This is a direct practical benefit of semantic space identity in
the first U(1) implementation.

\section{Operation Vocabulary}

\begin{center}
\begin{tabular}{ll}
\toprule
Operation & Contract \\
\midrule
\texttt{tensor\_product(A,B)}
  & ordered monoidal product \\
\texttt{compose(B,A)}
  & exact morphism composition \\
\texttt{repartition(A,...)}
  & boundary bend with duality toggle \\
\texttt{contract(A,B,pairs)}
  & checked partial-contraction planner \\
\texttt{trace(A,pair)}
  & categorical trace \\
\texttt{braid(A,...)}
  & explicit strand exchanges \\
\texttt{relabel(A,...)}
  & diagnostic metadata only \\
\bottomrule
\end{tabular}
\end{center}

The planner may lower these operations to axis permutations, dense
GEMMs, elementwise kernels, or communication.  Those lower operations
do not decide space compatibility, fusion legality, or planarity.

\section{Initial U(1) Implementation Sequence}

The first host-only U(1) \(\BlockTensor\) does not require general
\(F\)- or \(R\)-symbol machinery.  It should nevertheless establish
the following sequence:

\begin{enumerate}
  \item Add immutable identity-bearing spaces over the implemented
        structural decompositions.
  \item Represent ordered domain and codomain boundaries.
  \item Derive block legality from the all-out boundary.
  \item Keep the block key extensible with an empty initial coupling
        descriptor.
  \item Check exact space identity in composition and contraction.
  \item Implement the cheap U(1) form of repartition.
  \item Make SVD and truncation create and share one new bond space.
  \item Carry boundary and space metadata through async worklists,
        reduced-block kernels, CUDA lowering, and later MPI placement.
\end{enumerate}

There is no implicit dense fallback for a symmetry-typed operation.
An explicitly named dense projection may exist for diagnostics, but it
leaves the U(1) execution path and cannot feed data back into the
symmetry-aware state.

\section{Interface Invariants}

The following invariants refine the planar-network invariants and
should remain testable:

\begin{enumerate}
  \item Every tensor has an ordered domain and ordered codomain.
  \item Object duality is independent of morphism side.
  \item A dual object retains its primal space identity.
  \item Repartition toggles explicit duality and respects order
        reversal.
  \item Composition checks space identity and boundary compatibility.
  \item Isomorphic but distinct spaces require an explicit
        isomorphism.
  \item Labels and diagram coordinates never change legality.
  \item Block selection is an invariant-channel query over the
        oriented boundary.
  \item Global planarity belongs to the network or planner.
  \item Crossings, recouplings, and traces are explicit categorical
        operations.
  \item Storage movement preserves all boundary, sector, and coupling
        metadata.
\end{enumerate}

\subsection{Required Tests}

The initial test suite should include:

\begin{itemize}
  \item two isomorphic equal-dimensional spaces with different
        identities that cannot be composed accidentally;
  \item U(1) MPS and MPO block selection using the signed boundary
        equations above;
  \item round-trip repartition returning the original boundary and
        data;
  \item multiple-leg repartition with the required order reversal;
  \item SVD factors sharing exactly one new bond identity;
  \item relabeling that has no effect on contraction legality;
  \item worklist generation that retains every logical block key and
        space identity.
\end{itemize}

\section{Conclusion}

The local tensor boundary has two independent categorical axes:
\(\Space/\DualSpace\) and \(\Domain/\Codomain\).  Treating them as one
``direction'' type obscures wire bending, makes general contraction
harder to state, and loses information needed by non-Abelian and
braided categories.

The practical rule is:

\begin{quote}
Select blocks from an oriented morphism boundary, compose through
identity-bearing spaces, bend wires explicitly, and leave global
planarity to the contraction planner.
\end{quote}

For bosonic U(1), these rules reduce to cheap identity checks and signed
charge bookkeeping.  They are therefore suitable for the first
host-only DMRG implementation while preserving the structure required
for CUDA, MPI, non-Abelian fusion, and braiding.

\end{document}
